% !TEX root = ../../main_power.tex
% ---------------------------------------------------------------------
%  FreeRTOS und Low Power (Teil 6)
% ---------------------------------------------------------------------
\section{FreeRTOS und Low Power}
\label{sec:rtospwr}

\SW{Ausgangslage} FreeRTOS hat immer einen \T{Idle Task} mit der
niedrigsten Priorität. Er läuft, wenn kein anderer Task bereit ist --
also genau dann, wenn der Prozessor schlafen könnte, stattdessen aber in
einer Endlosschleife Strom verbrennt.

\subsection{Idle Hook}
\label{rez:idlehook}
Der Idle Task ruft optional bei jedem Durchlauf eine eigene Funktion auf.
Dort gehört der Übergang in den Low-Power-Modus hin.
\begin{lstlisting}
#define configUSE_IDLE_HOOK      1
#define configUSE_IDLE_HOOK_NAME McuRTOS_vApplicationIdleHook

void McuRTOS_vApplicationIdleHook(void) {
  __asm volatile("dsb");  /* Speicherzugriffe abschliessen */
  __asm volatile("wfi");  /* Wait For Interrupt -> schlafen */
  __asm volatile("isb");  /* Pipeline leeren */
}
\end{lstlisting}
\merkzeile Die Reihenfolge \R{dsb} -- \R{wfi} -- \R{isb} ist kein
Zufall: erst alle ausstehenden Speicherzugriffe abschliessen, dann
schlafen legen, nach dem Aufwachen die Pipeline neu füllen.

\begin{falle}[Der Tick weckt den Schläfer wieder]
Der \T{SysTick} läuft weiter und unterbricht den Idle Task bei jedem
Tick. Gemessen bei 1\,ms Tick: der Prozessor schläft jeweils nur
\T{978\,\textmu s}, wacht aber \T{1000-mal pro Sekunde} auf -- und jedes
Aufwachen kostet Zeit und Energie. Der Idle Hook allein holt deshalb nur
einen Teil des möglichen Gewinns.
\end{falle}

\subsection{Tickless Idle}
\label{rez:tickless}
FreeRTOS schaltet den Tick-Interrupt für die Dauer der erwarteten
Leerlaufzeit ganz ab und korrigiert den Tick-Zähler nach dem Aufwachen.
\begin{lstlisting}
#define configUSE_TICKLESS_IDLE  1

void McuRTOS_vApplicationIdleHook(void) {
  #if !configUSE_TICKLESS_IDLE   /* sonst doppelt */
  __asm volatile("dsb"); __asm volatile("wfi"); __asm volatile("isb");
  #endif
}

void McuRTOS_vOnPreSleepProcessing(TickType_t expectedIdleTicks) {
  __asm volatile("dsb");
  __asm volatile("wfi");
  __asm volatile("isb");
}
\end{lstlisting}
\merkzeile Ergebnis in der Trace-Ansicht: der Idle Task wird
\emph{nicht mehr unterbrochen} -- ein einziger langer Schlaf statt
tausend kurzer. \R{expectedIdleTicks} sagt, wie lange geschlafen werden
darf; danach muss ein Timer rechtzeitig wecken.

\begin{ztab}{@{}lX@{}}
\multicolumn{2}{@{}l}{\textbf{Die drei Stufen, aufsteigend sparsam}}\\
nichts & Idle Task dreht leer, voller \R{RUN}-Verbrauch \\
Idle Hook & schläft zwischen den Ticks, wacht $1/t_{tick}$-mal pro
  Sekunde \\
Tickless Idle & schläft über viele Ticks hinweg am Stück \\
Scheduler beenden & letzte Möglichkeit, wenn gar nichts mehr laufen
  soll \\
\end{ztab}
